\section{The finite evolution and its total force}\label{sec:finite}

Let $\cP_r$ be the space of positive Hermitian $r\times r$ matrices
of determinant one. Its affine-invariant metric is
\[
 \langle V,W\rangle_H=\tr(H^{-1}VH^{-1}W),
 \qquad
 d(H,K)=\|\log(H^{-1/2}KH^{-1/2})\|_{\HS}.
\]
It is a finite-dimensional Hadamard manifold. On $\cP_r^4$ use the
product metric $g_A=\sum_i A_i g_i$ and write $o=(I,I,I,I)$.
The masses are $A_i=1$ in our family; retaining them in the formulas
indicates the origin of each coefficient. All matrix norms below are
Hilbert--Schmidt norms unless stated otherwise.

For a Hermitian matrix $T$ set
\begin{equation}\label{eq:component-potential}
 f_T(H)=\tr(TH^{-1}TH)-\tr(T^2).
\end{equation}
On the family of Section~\ref{sec:geometry} define
\begin{equation}\label{eq:potential}
 \Phi(H)=\sum_{i=+,-}\alpha_i^2\kappa_i f_{T_i}(H_i)
       +\pi\sum_{e=ij}\epsilon_e d(H_i,U_{ij}H_jU_{ij}^*)^2.
\end{equation}
The negative gradient equation is
\begin{equation}\label{eq:finite-flow}
 H_i^{-1}\dot H_i
 =-\frac{\alpha_i^2\kappa_i}{A_i}
       [T_i,H_i^{-1}T_iH_i]
  +\frac{2\pi}{A_i}\sum_{e=ij}\epsilon_e B_{ij},
 \qquad B_{ij}=\log(H_i^{-1}U_{ij}H_jU_{ij}^*).
\end{equation}
The component term is omitted at rational vertices. The logarithm in
$B_{ij}$ is the principal matrix logarithm; it is $H_i$-self-adjoint
and trace-free. The edge summation uses the orientation from the vertex
under consideration. To check the component term, vary $H$ by $HX$,
with $X$ $H$-self-adjoint and trace-free. Differentiating
\eqref{eq:component-potential} gives
\[
 df_T(H)[HX]=\tr\bigl([T,H^{-1}TH]X\bigr).
\]
The gradient of half the squared distance to $K$ is
$-H\log(H^{-1}K)$, giving the edge term. The later reconstruction
derives the same coefficients from the Chern equation on the compact
fibre.

The following finite estimate is useful because it has no dependence
on a smallest positive weighted eigenvalue. We prove it on an arbitrary
fixed graph, allowing several component matrices at a vertex.

\begin{theorem}\label{thm:finite-force}
Fix a finite graph with $N$ vertices, rank $r\ge2$, positive masses
$A_i$, unitary edge matrices $U_e$, and a finite list of Hermitian
matrices $T_a$ at vertices $i(a)$. For arbitrary constant nonnegative
weights $w_e,\beta_a$ put
\begin{align}
 \Phi_{w,\beta}(H)&=\frac12\sum_e w_e
       d(H_i,U_eH_jU_e^*)^2+\sum_a\beta_a f_{T_a}(H_{i(a)}),
       \label{eq:general-potential}\\
 \cF_{w,\beta}(H)&=\sum_e w_e d(H_i,U_eH_jU_e^*)
           +\sum_a\beta_a\sqrt{2f_{T_a}(H_{i(a)})}.
       \label{eq:force}
\end{align}
There are $R_*>0$ and $c,C>0$, depending only on the fixed data,
such that, on $d_A(H,o)\le R_*$,
\begin{equation}\label{eq:force-gradient}
 c\cF_{w,\beta}(H)\le
 \|\grad\Phi_{w,\beta}(H)\|_{g_A}
 \le C\cF_{w,\beta}(H).
\end{equation}
The negative gradient trajectory starting at distance $R\le R_*$
exists for all $s\ge0$, stays in that ball, and satisfies
\begin{equation}\label{eq:force-integral}
 \int_0^\infty\bigl(\cF_{w,\beta}(H(s))
                   +\|\dot H(s)\|_{g_A}\bigr)\,ds\le CR,
 \qquad
 \int_0^\infty\Phi_{w,\beta}(H(s))\,ds\le\frac12R^2.
\end{equation}
It converges to a minimizer. These constants are independent of all
weights, their ratios and the subset of weights that vanish.
\end{theorem}

\subsection{A weighted linear estimate}
The elementary ingredient is a formula for a weighted least-squares
left inverse. Its constant retains the geometry of the fixed rows.

\begin{lemma}\label{lem:weighted-matrix}
For every fixed real matrix $D$ there is $C_D<\infty$ such that
\begin{equation}\label{eq:weighted-matrix}
 \|WDx\|\le C_D\|D^TWDx\|
\end{equation}
for every nonnegative diagonal matrix $W$ and every $x$.
\end{lemma}

\begin{proof}
First suppose $D$ is an $m\times n$ matrix of full column rank and
$W$ is positive. Let $D_S=R_SD$ be its restriction to $n$ rows, and
let $\mathfrak B$ consist of the subsets for which $\det D_S\ne0$.
Set $G=D^TWD$. Cauchy--Binet gives
\[
 \det G=\sum_{|S|=n}\det(D_S)^2\prod_{j\in S}w_j.
\]
In particular the numbers
\[
 p_S=\frac{\det(D_S)^2\prod_{j\in S}w_j}{\det G}
 \quad(S\in\mathfrak B)
\]
are nonnegative and sum to one. Replacing column $k$ of $D$ by $y$
replaces column $k$ of $D^TWD$ by $D^TWy$. Cauchy--Binet for this
replacement, followed by Cramer's rule, gives
\begin{equation}\label{eq:weighted-left-inverse}
 G^{-1}D^TW=\sum_{S\in\mathfrak B}p_S D_S^{-1}R_S.
\end{equation}
Terms with singular $D_S$ already vanish before Cramer's rule is
applied. Each $R_S$ has norm one, so the norm of
\eqref{eq:weighted-left-inverse} is at most
$\max_{S\in\mathfrak B}\|D_S^{-1}\|$. Transposing and applying
the result to $Gx$ proves \eqref{eq:weighted-matrix}.

For a matrix of rank $p>0$, choose an isometry $Q:\RR^p\to\RR^n$
onto $(\ker D)^\perp$. Apply the preceding argument to $D'=DQ$.
The identities $D=D'Q^T$ and $D^T=Q(D')^T$ preserve the norm on
the right of \eqref{eq:weighted-matrix}. Finally replace $W$ by
$W+\delta I$ and let $\delta\downarrow0$. This limit uses the
inequality, not an inverse of a singular normal matrix. The case $D=0$
is immediate.
\end{proof}

If $D$ is divided into $J$ fixed blocks $D_k$ and $W$ has constant
weight $\lambda_k\ge0$ on each block, the lemma gives
\begin{equation}\label{eq:linear-force}
 F^0(x):=\sum_k\lambda_k\|D_kx\|
 \le\sqrt J\,C_D\|D^TWDx\|.
\end{equation}
The choice of $(\ker D)^\perp$ is made once for the full fixed
matrix. It does not change with the support of $W$.

\subsection{Exact zero sets and nonlinear absorption}
Use real coordinates $x_i=\sqrt{A_i}\log H_i$. Then
$|x|=d_A(H,o)$ and the pullback metric is Euclidean at zero. The
linear edge maps and component maps are
\[
 D_ex=\frac{x_i}{\sqrt{A_i}}-
       \Ad(U_e)\frac{x_j}{\sqrt{A_j}},
 \qquad
 C_ax=i[T_a,x_{i(a)}/\sqrt{A_{i(a)}}].
\]
Both maps take values in real spaces of Hermitian matrices. Associate
to an edge the function $f_e=\frac12d(H_i,U_eH_jU_e^*)^2$, and
to a component matrix the function $f_a=f_{T_a}(H_{i(a)})$.
The exact zero set of $f_e$ is $\ker D_e$, and its Hessian at zero
is $D_e^TD_e$.

For the component assertion, diagonalize $S=\log H$ with eigenvalues
$s_b$. In this basis $T$ is Hermitian and direct multiplication gives
\begin{equation}\label{eq:cosh-potential}
 f_T(e^S)=2\sum_{b<c}|T_{bc}|^2
            \bigl(\cosh(s_b-s_c)-1\bigr)
 \ge\frac12\|[T,S]\|^2.
\end{equation}
Its exact zero set is $[T,S]=0$, and its quadratic term is
$\frac12\|[T,S]\|^2$. Thus the component Hessian is $C_a^TC_a$.
Stack all these maps into $D$. Denote any one block by $D_k$ and
its associated function by $f_k$. Every $f_k$ is nonnegative,
vanishes precisely on $\ker D_k$, and has Hessian $D_k^TD_k$ at
zero. In particular its gradient vanishes on the whole kernel.

On the auxiliary ball $|x|\le1$ the third derivatives of this finite
collection of functions are bounded. Put
$r_k(x)=\nabla f_k(x)-D_k^TD_kx$. On $|x|\le R\le1$,
$\|Dr_k(x)\|\le C_0R$. If $P_k$ is orthogonal projection onto
$\ker D_k$, then $r_k(P_kx)=0$ and the segment from $P_kx$ to
$x$ stays in the ball. Therefore
\begin{equation}\label{eq:block-remainder}
 \|r_k(x)\|\le C_0R\dist(x,\ker D_k)
                \le C_1R\|D_kx\|.
\end{equation}
The last constant uses only the positive singular values of each
nonzero fixed block. Zero blocks can be omitted. Combining
\eqref{eq:block-remainder} with \eqref{eq:linear-force}, and writing
$K=\sqrt J C_D$, gives
\[
 F^0(x)\le K\|\nabla\Phi_{w,\beta}(x)\|
                    +KC_1R F^0(x).
\]
Choose $R_*$ so that $KC_1R_*\le\frac12$. Absorption gives
$F^0\le2K\|\nabla\Phi_{w,\beta}\|$ uniformly in the weights.

For edges, the distance-decreasing logarithm at the identity and the
bounded differential of the exponential on this ball imply
$\|D_ex\|\le\sqrt{2f_e(x)}\le C\|D_ex\|$.
For components, the lower bound is \eqref{eq:cosh-potential}; Taylor
expansion along the segment from $P_ax$ to $x$, on which the initial
function and gradient vanish, gives the upper bound. Hence
$F^0\le\cF_{w,\beta}\le CF^0$. Smooth equivalence of the
coordinate metric with the Euclidean metric proves the left side of
\eqref{eq:force-gradient}. The block bounds
$\|\nabla f_k\|\le C\|D_kx\|$ prove the right side.

\subsection{Convexity and finite length}
To verify geodesic convexity of $f_T$, write
$H(t)=H_0^{1/2}e^{tK}H_0^{1/2}$ and diagonalize the Hermitian matrix
$K$ with eigenvalues $k_b$. For $B=H_0^{-1/2}TH_0^{1/2}$,
\[
 \tr(TH(t)^{-1}TH(t))
       =\sum_{b,c}|B_{cb}|^2 e^{t(k_b-k_c)}.
\]
Each summand is convex. Squared distance between two moving geodesics
is convex on a Hadamard manifold, and unitary conjugation is an
isometry. Thus \eqref{eq:general-potential} is geodesically convex,
nonnegative and minimized at $o$. Its negative gradient flow obeys
\begin{equation}\label{eq:finite-energy}
 \frac12\frac{d}{ds}d_A(H(s),o)^2\le-\Phi_{w,\beta}(H(s)).
\end{equation}
It stays in the compact ball of radius $R$ and exists for all time.
Integration gives the second estimate of \eqref{eq:force-integral}.

For $T>s$, convexity and decreasing potential similarly give
\[
 \frac12\frac{d}{ds}d_A(H(s),H(T))^2
       \le\Phi_{w,\beta}(H(T))-\Phi_{w,\beta}(H(s))\le0.
\]
The curve is therefore self-contracted: its distance to any later
point decreases as its earlier time increases. We apply Ohta's
bounded-image length estimate on a finite-dimensional Hadamard
manifold \cite[Theorem~3.3 and (3.8)]{ohta}. Here the dimension is
$N(r^2-1)\ge3$, the image has diameter at most $2R$, and its
radius-one neighborhood lies in $B(o,R_*+1)$, whose volume is fixed.
The length estimate consequently gives
$\int_0^\infty\|\dot H\|_{g_A}\,ds\le CR$.
Together with \eqref{eq:force-gradient}, this proves the first
estimate of \eqref{eq:force-integral}.

Finite length implies convergence to a point $H_\infty$ in the ball.
If its gradient were nonzero, continuity would give a positive lower
bound for the speed for all sufficiently large times, a contradiction.
Convexity makes $H_\infty$ a global minimizer, so all terms with
positive weights vanish there. In logarithm coordinates the minimizer
set is exactly $\bigcap_{\lambda_k>0}\ker D_k$. The limit need
not be $o$, and the theorem gives neither a uniform rate nor an
explicit formula for this selected limit. This completes the proof
of Theorem~\ref{thm:finite-force}.

\subsection{Dependence on the fixed constraints}
The fixed-matrix hypothesis cannot be replaced by uniform injectivity
of the full matrix. For $\delta>0$ take
\[
 D_\delta=\begin{pmatrix}1&0\\0&1\\1&\delta\end{pmatrix},
 \quad W_\delta=\operatorname{diag}(1,\delta^4,1),
 \quad x_\delta=\begin{pmatrix}-\delta/2\\1\end{pmatrix}.
\]
Although $D_\delta^TD_\delta\ge I$,
\[
 W_\delta D_\delta x_\delta
 =(-\delta/2,\delta^4,\delta/2)^T,
 \qquad D_\delta^TW_\delta D_\delta x_\delta
 =(0,\delta^2/2+\delta^4)^T.
\]
Their norm ratio is asymptotic to $\sqrt2/\delta$.
A row basis that can dominate the weighted normal matrix becomes
singular. This is an obstruction to a uniform force estimate, not
an obstruction theorem for actual HYM approximation.

There is a second distinction needed in the analytic proof. A
component force is weighted by $\alpha^2$, whereas a change of
component frame can have a first-order factor $\alpha$. With one
vertex, mass one, $T=\sigma_z$, $\kappa=1$ and
$H=\exp(q\sigma_x)$, direct calculation gives
\[
 f_T(H)=4\sinh^2q,\qquad
 g(\partial_qH,\partial_qH)=2,\qquad
 \dot q=-2\alpha^2\sinh(2q).
\]
For $q(0)=q_0>0$ and $\alpha>0$ it follows that
\begin{equation}\label{eq:first-order-divergence}
 \int_0^\infty\alpha q(s)\,ds
   =\frac1{2\alpha}\int_0^{q_0}\frac{q}{\sinh(2q)}\,dq.
\end{equation}
The right-hand integral is finite and strictly positive. Thus a
first-order frame error cannot be integrated uniformly by
Theorem~\ref{thm:finite-force} alone. The angular divergence
construction below addresses precisely this term.
